%
% reversi_140_corpus_analisi.tex
%
% This file is part of the reversi program
% http://github.com/rcrr/reversi-v2
%
% Author Roberto Corradini mailto:rob_corradini@yahoo.it
% Copyright 2026 Roberto Corradini. All rights reserved.
%
% License
%
% This program is free software; you can redistribute it and/or modify it
% under the terms of the GNU General Public License as published by the
% Free Software Foundation; either version 3, or (at your option) any
% later version.
%
% This program is distributed in the hope that it will be useful,
% but WITHOUT ANY WARRANTY; without even the implied warranty of
% MERCHANTABILITY or FITNESS FOR A PARTICULAR PURPOSE. See the
% GNU General Public License for more details.
%
% You should have received a copy of the GNU General Public License
% along with this program; if not, write to the Free Software
% Foundation, Inc., 59 Temple Place - Suite 330, Boston, MA  02111-1307, USA
% or visit the site <http://www.gnu.org/licenses/>.
%

%-----------------------------------------------------------------------------
% Capitolo 8: Corpus e analisi delle feature.
%
% La fase della pipeline che trasforma il corpus grezzo di partite in dati di
% feature indicizzati e analizzati: i livelli SU_ANALYSIS, POSITIONS, INDEXES,
% ANOVA e WMAPS.
%-----------------------------------------------------------------------------

\chapter{Corpus e analisi delle feature}
\label{ch:corpus-analisi}

Questo capitolo copre la \emph{seconda} delle cinque fasi della pipeline —
\textbf{Corpus e analisi delle feature} — così come il capitolo sulla pipeline
la raggruppa (tabella~\ref{tab:pipeline-levels}). La fase è formata da cinque
livelli, eseguiti in quest'ordine: \texttt{SU\_ANALYSIS}, \texttt{POSITIONS},
\texttt{INDEXES}, \texttt{ANOVA} e \texttt{WMAPS}. Il suo compito è
\emph{trasformare il corpus grezzo di partite nei dati di feature
indicizzati e analizzati} che il costruttore del modello (la fase successiva,
\emph{Costruzione del modello}) ha bisogno.

Il punto di partenza della fase è il \texttt{feature\_set} già risolto dalla
fase Setup: l'insieme ordinato delle feature (l'intercetto, le mobility, i
pattern) con, per ognuna, il numero di istanze e lo spazio delle
configurazioni. Il punto di arrivo è un pacchetto compatto di artefatti: le
posizioni di addestramento e i loro valori di gioco, la matrice degli indici
categorici, le mappe che collegano ogni configurazione di feature a un indice
del vettore di pesi. Tutto il resto che la fase produce — le tabelle di
frequenza, le misure di varianza, i report di rarità — è \emph{diagnostica}:
informa il progettista sul \emph{perché} un certo feature set è buono o
cattivo, ma non entra nel modello.

C'è una distinzione di peso che conviene fissare subito. Dei cinque livelli,
tre sono \emph{obbligatori} e formano la spina dorsale della fase:
\texttt{POSITIONS}, \texttt{INDEXES} e \texttt{WMAPS}. Senza di essi non c'è
nessun modello. Gli altri due — \texttt{SU\_ANALYSIS} e \texttt{ANOVA} — sono
\emph{opzionali}: la loro flag \texttt{enable} è \texttt{false} di default,
quindi una run standard non li esegue. Quando sono attivi, però, producono
l'analisi più profonda che lo strumento offre sulla qualità delle feature. E
c'è un accorgimento di \emph{degradazione elegante}: l'\texttt{ANOVA}, se
trova l'\texttt{SU\_ANALYSIS} spento (quindi nessun corpus nel contesto), non
fallisce — salta semplicemente la parte che del corpus avrebbe bisogno e
prosegue con il resto.

La tabella riassume i cinque livelli: da cosa legge e cosa pubblica sul
contesto.

\begin{table}[htbp]
\centering
\small
\begin{tabularx}{\textwidth}{@{}llXX@{}}
\toprule
Livello & Worker & Legge dal contesto & Pubblica sul contesto \\
\midrule
\texttt{SU\_ANALYSIS} & \texttt{wrk\_su\_analysis} & \texttt{cfg.su\_analysis}, \texttt{feature\_set}, \texttt{empty\_count} & \texttt{su\_analysis\_tables} \\
\texttt{POSITIONS}    & \texttt{wrk\_positions}     & \texttt{cfg.regab\_data\_set}, \texttt{empty\_count} & \texttt{positions}, \texttt{game\_values}, \texttt{rds\_checksum}, \texttt{training\_pop\_stats} \\
\texttt{INDEXES}      & \texttt{wrk\_indexes}       & \texttt{positions}, \texttt{feature\_set}, \texttt{rds\_checksum} & \texttt{hlp\_indexes}, \texttt{hlp\_indexes\_checksum} \\
\texttt{ANOVA}        & \texttt{wrk\_anova}         & \texttt{hlp\_indexes}, \texttt{game\_values}, \texttt{su\_analysis\_tables}, \texttt{cfg.anova} & (solo i file di report) \\
\texttt{WMAPS}        & \texttt{wrk\_wmaps}         & \texttt{hlp\_indexes}, \texttt{feature\_set}, \texttt{cfg.stat\_model} & \texttt{feature\_w\_ranges}, \texttt{iwmap\_feature\_offset}, \texttt{iwmap}, \texttt{wmap}, \texttt{wmap\_fallback}, \texttt{w}, \texttt{wmaps\_obj\_checksum} \\
\bottomrule
\end{tabularx}
\caption{I cinque livelli della fase: la lettura e la pubblicazione sul contesto condiviso.}
\label{tab:corpus-levels}
\end{table}

Prima di entrare nei singoli livelli, una nota sul meccanismo che quattro di
loro condividono. \texttt{SU\_ANALYSIS} (la sola parte corpus),
\texttt{POSITIONS}, \texttt{INDEXES} e \texttt{WMAPS} usano tutti lo stesso
\emph{flusso di cache a quattro passi}, già introdotto nel capitolo sulla
pipeline: a ogni run il worker affida al gestore della cache l'artefatto del
livello, che (1) \emph{cerca} il file cacheato, (2) ne \emph{valida} la firma
contro la configurazione live, (3) in caso di file assente o firma
inconsistente lo \emph{ricalcola} e lo \emph{riporta} su file, (4) in caso
contrario lo \emph{restituisce} come già pronto. La validazione è sempre un
\emph{test puro} — nessun database, nessun ricalcolo: un confronto tra pochi
campi, e il dettaglio di \emph{quali} campi è l'oggetto delle singole
sezioni che seguono. Due eccezioni da tenere a mente: l'\texttt{ANOVA} non ha
artefatto cacheabile (scrive solo report, e li riscrive a ogni run), e
l'\texttt{SU\_ANALYSIS} applica il flusso alla sola metà \emph{costosa} del
suo lavoro — la generazione del corpus — mentre le statistiche, economiche,
sono ricalcolate sempre.

\section{Il livello \texttt{SU\_ANALYSIS}}
\label{sec:corpus-su}

\texttt{SU\_ANALYSIS} è il livello più ricco della fase e, nella run, il
più costoso: è quello che \emph{genera} un corpus di posizioni per studiarne
le feature. Prima di entrare nel meccanismo, va chiarito cosa si intende per
\emph{analisi SU}.

\subsection{Cos'è l'analisi SU}

La sigla \texttt{SU} sta per \emph{Symmetrical Uncertainty},
\emph{incertezza simmetrica}. È una misura della \emph{dipendenza} tra due
variabili categoriche, costruita a partire dall'informazione mutua e
\emph{normalizzata} per compensare il bias matematico intrinseco
dell'informazione mutua grezza (che tende a sovrastimare variabili con molte
categorie). Se $u$ e $v$ sono due variabili categoriche, con entropie
marginali $H(u)$ e $H(v)$ e informazione mutua $I(u;v)$, l'incertezza
simmetrica è

\begin{equation}
  \mathrm{SU}(u,v) \;=\; \frac{2\,I(u;v)}{H(u) + H(v)} \;\in\; [0,1].
  \label{eq:su}
\end{equation}

Il valore è \emph{simmetrico} (invariare $u$ e $v$ non lo cambia) e
\emph{normalizzato}: $0$ significa ortogonalità perfetta (le due variabili
portano informazioni indipendenti), $1$ significa che le due variabili sono
la \emph{stessa} informazione. Il report allega una legenda fissa per
interpretare i valori:

\begin{center}
\small
\begin{tabular}{ll}
\toprule
Intervallo di \texttt{SU} & Interpretazione \\
\midrule
$0.0$            & ortogonalità perfetta \\
$0.0 - 0.1$      & eccellente indipendenza \\
$0.1 - 0.3$      & correlazione moderata \\
$0.3 - 0.5$      & correlazione forte \\
$> 0.5$          & correlazione molto forte \\
\bottomrule
\end{tabular}
\end{center}

Perché questa misura interessa un costruttore di modelli? Perché le feature di
un buon set di valutazione dovrebbero essere, a coppie, \emph{quasi
ortogonali}: se due feature sono fortemente correlate, una delle due è
redundante e aggiunge parametri al modello senza aggiungere informazione
predittiva. L'analisi SU risponde alla domanda ``quale feature duplica
quale?''.

L'unità di analisi non è la feature intera, ma la \emph{sotto-variabile}:
la coppia (feature, istanza), che corrisponde a \emph{una colonna} della
matrice degli indici. Ogni istanza di un pattern è una sua variabile
categorica a sé; una mobility è una singola colonna. È la stessa unità su cui
vivono i pesi del modello, e su cui l'\texttt{ANOVA} (più avanti) misurerà la
varianza spiegata.

\subsection{La generazione del corpus}

A differenza dei livelli a valle, che leggono le posizioni dal database,
\texttt{SU\_ANALYSIS} \emph{genera} le proprie posizioni da zero, giocando
partite casuali. La ragione è metodologica: per studiare la \emph{frequenza
naturale} con cui le configurazioni di feature si presentano nel gioco, serve
un campione \emph{neutro}, non condizionato dalla selezione che il database
ha già applicato. Il corpus è quindi una popolazione di posizioni casuali,
tutte \emph{arrestate} alla stessa fase di gioco.

Lo spec del generatore vive nella sezione \texttt{su\_analysis.corpus} della
configurazione: il \texttt{seed} (default 42) del generatore di numeri
casuali MT19937, e il \texttt{total\_positions} (default 10\,000\,000) di
posizioni da raccogliere. La regola di arresto — fino a quale numero di
caselle vuote portare ogni partita — è l'\texttt{empty\_count} \emph{di primo
livello}, in modo che il corpus si fermi alla \emph{stessa} fase di gioco
usata dall'estrazione dal database per l'addestramento.

La generazione è \emph{a onde} (\emph{wave}) e \emph{in streaming}. Si parte
dalla posizione iniziale canonica del gioco; a ogni onda si lanciano
$n$ \emph{lane} in parallelo (tutte copie della posizione iniziale), dove
$n$ è il minimo tra una costante di onde (3\,000\,000) e la quota ancora
mancante; ogni lane gioca mosse casuali legali finché non raggiunge
l'\texttt{ec} target oppure la partita termina. Al ritorno, si \emph{raccoglie}
solo le lane che si sono fermate \emph{esattamente} al target: le lane la cui
partita è finita prima (con più caselle vuote del target) sono scartate, e
la quota viene integrata con l'onda successiva. Il tratto decisivo per la
memoria è che \emph{il corpus non viene mai materializzato} in RAM: a ogni
onda si calcola la matrice degli indici \emph{solo} delle lane raccolte e si
passa quella matrice agli accumulatori, che aggiornano i propri contatori.
L'onda si scarta e si passa alla successiva.

\begin{lstlisting}[language=python, caption={Il cuore di \texttt{hlp\_su\_analysis\_compute}: la generazione a onde del corpus, in streaming.}]
prng  = MT19937(corpus.seed)
igp   = initial_game_position()
marg  = _MarginalAccumulator(fset)
joint = _JointAccumulator(fset, su_pairs) if su_pairs else None

harvested = 0
while harvested < total:
    n     = min(WAVE_SIZE, total - harvested)
    batch = np.full(n, igp)                        # n copie della posizione iniziale
    final, ec_arr = random_game_play_path(batch, ec, prng)
    keep  = ec_arr == ec                           # solo le lane ferme a ec esatto
    idx   = fset.compute_indexes(final[keep])      # matrice degli indici dell'onda
    marg.add(idx)
    if joint is not None:
        joint.add(idx)
    harvested += int(keep.sum())
\end{lstlisting}

Il corpus così generato è \emph{deterministico} dato il seed: la stessa
configurazione riproduce lo stesso campione, posizione per posizione.

\subsection{Le tabelle marginali}

Il primo accumulatore, \texttt{\_MarginalAccumulator}, produce la
\emph{tabella marginale}: per ogni sotto-variabile (feature, istanza),
l'istogramma delle frequenze con cui ogni sua configurazione si è presentata
nel corpus. Lo schema di ogni riga è
$(\text{name}, \text{instance}, n_{\text{conf}}, \text{codes},
\text{counts})$: il nome della feature, l'indice dell'istanza, la dimensione
teorica dello spazio delle configurazioni, e — per il \emph{supporto
osservato} — i codici visti e i relativi conteggi.

C'è qui una scelta di memoria degna di nota. Lo spazio delle configurazioni di
una sotto-variabile può essere piccolo (una mobility a maschera piena ha
$65$ valori, da $0$ a $64$ mosse) o enorme (un pattern a $16$ caselle ne ha
$3^{16} \approx 43$ milioni). L'accumulatore usa una \emph{soglia}: se la dimensione teorica
è sotto un milione di codici, tiene un istogramma \emph{denso} (un vettore di
contatori, aggiornato con un \texttt{bincount} vettorizzato); oltre, passa a
un dizionario \emph{sparse} che registra solo i codici effettivamente visti,
fusi a ogni onda. Nel caso attuale tutte le sotto-variabili cadono nel ramo
denso (il massimo in gioco è $59\,049$), ma la via sparse è lì per i spazi di
codice più ampi.

\subsection{Le tabelle congiunte}

Il secondo accumulatore, \texttt{\_JointAccumulator}, è costruito \emph{solo}
se la configurazione dichiara delle coppie da analizzare
(\texttt{su\_pairs}, una lista di coppie ordinate di nomi di feature). Per
ogni coppia $(F_i, F_j)$ richiesta, e per ogni cella
$(\text{istanza}_i, \text{istanza}_j)$ del reticolo di istanze, l'accumulatore
tiene la \emph{tabella congiunta sparsa}: quante volte la configurazione $c_i$
della prima è apparsa \emph{insieme} alla configurazione $c_j$ della seconda.
Lo spazio $c_i \times c_j$ è troppo grande per essere denso, quindi si tiene
solo il supporto osservato, con la chiave dei due codici \emph{impaccata} in
un singolo intero a 64 bit ($\text{key} = c_i \cdot n_{\text{conf},j} + c_j$)
che viene poi disimpaccata in lettura.

Un accorgimento di correttezza: la cella \emph{diagonale} di una coppia
formata dalla \emph{stessa} feature (lo stesso nome, la stessa istanza,
$\text{istanza}_i = \text{istanza}_j$) \emph{non viene costruita}. La
congiunta di una variabile con se stessa \emph{è} la sua marginale —
calcolarla sarebbe ridondante, e la sua \texttt{SU} sarebbe bloccata a $1$
per costruzione, non una misura. Si analizzano invece le celle
\emph{fuori diagonale}, che per una stessa feature sono le sue \emph{varie
orientazioni $D_4$} tra loro: è la misura dell'\emph{anisotropia} di un pattern
(su quante delle sue istanze il gioco si concentra).

\subsection{Le statistiche: dispersione, rarità e incertezza simmetrica}

A differenza della generazione del corpus (costosa, cacheata), le
\emph{statistiche} sono economiche e vengono \emph{ricalcolate a ogni run},
anche quando le tabelle vengono lette dalla cache. Producono il report TXT
(più un CSV opzionale) in tre sezioni.

\emph{Sezione A — dispersione per sotto-variabile.} Per ogni sotto-variabile
il report calcola e tabula cinque grandezze:

\begin{itemize}
  \item $K_{\text{obs}}$: il numero di configurazioni \emph{distinte
        osservate} nel corpus;
  \item $H$: l'entropia di Shannon della distribuzione osservata, in bit,
        $H = -\sum_c p(c)\,\log_2 p(c)$;
  \item $H_{\text{norm}} = H / \log_2 K_{\text{obs}} \in [0,1]$: vale $1$ se
        l'osservato è uniforme sul supporto visto;
  \item \emph{coverage} $= K_{\text{obs}} / N_{\text{conf}}$: la frazione
        dello spazio teorico effettivamente coperta dal gioco casuale (parte
        delle configurazioni teoriche non è raggiungibile per le regole del
        gioco);
  \item \emph{never-obs} $= N_{\text{conf}} - K_{\text{obs}}$: le
        configurazioni teoriche \emph{mai} viste nel corpus.
\end{itemize}

Qui compare una distinzione sottile e importante tra i due tipi di feature.
Per una \emph{mobility}, i codici osservati sono i valori grezzi del conteggio
(nessuna riduzione per simmetria), quindi lo spazio teorico in gioco è
$N_{\text{conf}} = k+1$ (con $k$ il numero di caselle della maschera). Per un
\emph{pattern}, i codici osservati sono gli \emph{indici principali} — i
rappresentanti delle orbite $D_4$ — quindi lo spazio teorico non è il $3^{S}$
grezzo ma il numero $n_{\text{unique}}$ di \emph{classi di equivalenza}
(sotto le simmetrie indotte sulla maschera), dato dalla media di Burnside.
È questo il denominatore corretto per la coverage e per l'audit
never-observed di un pattern.

\emph{Sezione B — rarità (opzionale).} Se la configurazione dichiara
\texttt{rarity\_thresholds} (una lista strettamente crescente di soglie di
frequenza), il report aggiunge, per ogni sotto-variabile, un istogramma a
\emph{bin} delle frequenze delle singole configurazioni: bin semi-aperti
$[0, t_1)$, $[t_1, t_2)$, \ldots, $[t_k, +\infty)$. Il \emph{primo} bin
raccolge i codici a frequenza zero, cioè le configurazioni che il gioco
casuale \emph{quasi mai} raggiunge — e che, di conseguenza,
\emph{non potranno mai avere un peso appreso}. È la sezione che quantifica,
feature per feature, quanta parte del suo spazio di configurazioni è, ai
fatti, inerte.

\emph{Sezione C — incertezza simmetrica.} Per ogni coppia richiesta e ogni
cella fuori diagonale, il report calcola la \eqref{eq:su} a partire dalla
congiunta e dalle due marginali. Le formule intermedie sono
l'informazione mutua, calcolata per identità a partire dalle tre entropie,
\begin{equation}
  I(u;v) \;=\; H(u) + H(v) - H(u,v),
  \label{eq:mutual-info}
\end{equation}
dove $H(u,v)$ è l'entropia della congiunta (troncata a $\ge 0$ per
arrotondamento in virgola mobile), e poi la normalizzazione
\eqref{eq:su}. Essendo la misura simmetrica, le due direzioni $(i,j)$ e
$(j,i)$ della \emph{stessa} coppia sono la \emph{stessa} misura: il report le
\emph{deduplica} in un ordine canonico dei due estremi e le raggruppa per
coppia di feature, ordinando dentro ogni gruppo per ordine lessicografico
(feature, istanza).

Il report così composto ha una forma fissa; un estratto (con valori fittizi,
ma nel formato reale) ne mostra l'aspetto:

\begin{lstlisting}[caption={Estratto illustrativo del report \texttt{su\_analysis} (formato reale, valori fittizi).}]
Symmetrical Uncertainty Analysis Report
========================================

Corpus: 2,000,000 positions
Features: LMC, ALMC, EDGE, COR
SU pairs: EDGE - COR

Section A: Properties of features selected by the logistic model
K_obs    : distinct configurations observed (of the theoretical space)
H        : Shannon entropy of the observed distribution (bits)
H_norm   : H / log2(K_obs) in [0,1] - 1.0 = uniform over the observed
coverage : observed / theoretical - the ratio of observed to total
           theoretical configurations (some are unreachable by the game's rules)
never-obs: theoretical - observed (theoretical configurations never seen)
  LMC   inst 0: K_obs 65/65,    H 3.90,  H_norm 0.686, coverage 1.000, never-obs 0
  EDGE  inst 0: K_obs 4,211/4,366, H 8.71, H_norm 0.724, coverage 0.965, never-obs 155
  COR   inst 0: K_obs 3/3,      H 1.55,  H_norm 0.978, coverage 1.000, never-obs 0

Section B: Rarity report
The per-bin frequencies of the codes: the half-open bins [0, t1), [t1, t2), ..., [tk, inf)
The first bin the zero-frequency codes, the configurations random play
almost never reaches: the weights that can never be learned.
  LMC   inst 0: [0,1) 0 | [1,10) 2 | [10,100) 4 | [100,inf) 59
  EDGE  inst 0: [0,1) 155 | [1,10) 1,412 | [10,100) 1,180 | [100,inf) 1,619

Section C: Symmetrical uncertainty
It evaluates the normalized relationship, dependency, or shared variance
between variables, compensating for the inherent mathematical bias found in
raw information gain.
SU        : Meaning
0.0       : Perfect orthogonality
0.0 - 0.1 : Excellent independence
0.1 - 0.3 : Moderate correlation
0.3 - 0.5 : Strong correlation
over 0.5  : Very strong correlation
EDGE - COR
  SU EDGE[0] - COR[0] = 0.3127
  SU EDGE[1] - COR[0] = 0.2418
  SU EDGE[2] - COR[0] = 0.1975
\end{lstlisting}

\subsection{Cache e pubblicazione}

La parte costosa (il corpus) è cacheata con la firma di consistenza
\emph{più} ricca della fase: lo hash di dominio del \texttt{feature\_set}, il
\texttt{seed}, il \texttt{total\_positions}, l'\texttt{ec}, la versione del
generatore di corpus e la versione del PRNG, e la lista \texttt{su\_pairs}.
Il punto di questa firma è che \emph{ogni} cambio di comportamento del
generatore (ad esempio un aumento della dimensione di onda, che sposta il
flusso MT19937 e produce corpus \emph{diversi}) invalidi ogni corpus cacheato:
la firma esiste proprio perché un cambiamento di generatore \emph{non possa
mai} essere riusato in silenzio. A run conclusa, il livello pubblica
\texttt{ctx.su\_analysis\_tables} (il dizionario con firma, marginali e
congiunte) — che è proprio ciò di cui l'\texttt{ANOVA} si servirà come
\emph{distribuzione ideale} del corpus — e, in caso di disabilitazione, lo
porta a \texttt{None}.

\section{Il livello \texttt{POSITIONS}}
\label{sec:corpus-positions}

\texttt{POSITIONS} è il primo dei tre livelli obbligatori: carica le
\emph{posizioni di addestramento} e i loro \emph{valori di gioco veri}. A
differenza del corpus sintetico di \texttt{SU\_ANALYSIS}, qui i dati sono
\emph{reali}: provengono dal database REGAB, il deposito delle posizioni
giocate casualmente e poi \emph{risolte} (ciascuna taggata con il suo valore
vero di partita, il differenziale finale di dischi in $[-64, +64]$).

\subsection{La fonte: il database REGAB}

La sezione \texttt{regab\_data\_set} della configurazione descrive la
connessione al database (\texttt{dbname}, \texttt{user}, \texttt{host},
\texttt{port}, \texttt{password}) e i \emph{criteri di filtro} dell'estrazione:
la lista di \texttt{bid} (identificatori di torneo) e la lista di
\texttt{status} (codici di partita a tre lettere). L'estrazione SQL seleziona
dalla tabella delle posizioni le righe che appartengono a quei tornei, con
quello \texttt{empty\_count} e quegli status, ordinate per sequenza:

\begin{lstlisting}[language=sql, caption={La forma dell'estrazione dal database (l'istruzione costruita da \texttt{regab\_gp\_as\_df}).}]
SELECT mover, opponent, game_value
FROM   regab_prng_gp
WHERE  batch_id  IN (<bid>)
  AND  empty_count = <ec>
  AND  status    IN (<status>)
ORDER  BY seq;
\end{lstlisting}

Il risultato è un \texttt{RegabDataSet}: un oggetto che porta, oltre al
dataframe \texttt{df\_mogv} con le tre colonne \texttt{mover}, \texttt{opponent}
e \texttt{game\_value}, i metadati della query (\texttt{bid}, \texttt{status},
\texttt{ec}) e il riferimento alla connessione. Le colonne
\texttt{mover}/\texttt{opponent} sono interi a 64 bit, cioè \texttt{Bitboard};
\texttt{game\_value} è un intero a 8 bit (\texttt{int8}), il che riflette
proprio il suo range $[-64, +64]$.

\subsection{La politica sulle duplicazioni}

Nello stesso \emph{flusso di calcolo} (quello che popola la cache), subito
dopo l'estrazione, il livello applica la politica sulle
\emph{duplicazioni} governata dalla flag
\texttt{regab\_allow\_duplicate} di primo livello. Una posizione può essere
registrata nel database più di una volta (ad esempio sotto status diversi),
e la query di estrazione restituisce \emph{tutte} le occorrenze. La chiave di
deduplicazione è la coppia $(\text{mover}, \text{opponent})$: se
\texttt{regab\_allow\_duplicate} è \texttt{false}, le righe in eccesso sono
rimosse, tenendo la \emph{prima} occorrenza in ordine di estrazione (che,
essendo la query ordinata per \texttt{seq}, è quella di \texttt{seq} minimo).
Se è \texttt{true} (il default), ogni riga è conservata. In entrambi i casi il
conteggio delle duplicazioni è \emph{loggato}.

La ragione per cui la deduplicazione vive \emph{dentro} il flusso di calcolo
(non a valle) è la cache: il file cacheato contiene il data set
\emph{post}-politica, quindi cambiare la flag cambia il contenuto salvato e il
\texttt{rds\_checksum}, che invalida di conseguenza i livelli a valle lungo la
catena di checksum già esistente — senza alcun cambio di formato file.

\subsection{Le posizioni e i valori di gioco}

Fatto il carico (dalla cache o dal calcolo), il livello applica il
\emph{post-processing} e pubblica gli artefatti. Il metodo
\texttt{generate\_positions\_and\_game\_values()} del data set converte il
dataframe nelle due rappresentazioni che il resto della pipeline consuma:

\begin{itemize}
  \item \texttt{positions}: un array strutturato di \texttt{Position} —
        ogni record è la coppia $(\text{mover}, \text{opponent})$, ciascun
        campo un \texttt{Bitboard} (\texttt{uint64}, bit 0 = casella A1);
  \item \texttt{game\_values}: un array \texttt{int8} dei valori di gioco veri,
        allineato posizione per posizione.
\end{itemize}

Accanto a questi, il livello pubblica \texttt{rds\_checksum} (l'impronta
SHA3-256 dei byte grezzi del data set, che farà da chiave ai cache a valle) e
\texttt{training\_pop\_stats}, un dizionario con le statistiche
\emph{di popolazione} del set di addestramento: \texttt{count}, \texttt{mean},
\texttt{std\_dev} e \texttt{variance} dei valori di gioco. Queste ultime sono
il riferimento contro cui, più avanti, si misurerà la qualità predittiva del
modello: se la varianza della popolazione è $V$, una feature che ne spiega una
frazione significativa è, per definizione, una feature utile.

La consistenza della cache confronta i parametri di connessione
(\texttt{dbname}, \texttt{user}, \texttt{host}, \texttt{port}) \emph{e} i
parametri di query (\texttt{bid}, \texttt{status}, \texttt{ec}): se
\emph{qualcuno} è diverso tra la configurazione live e quella cacheata, la
cache è scartata e il data set è riestratto.

\section{Il livello \texttt{INDEXES}}
\label{sec:corpus-indexes}

\texttt{INDEXES} è il livello che \emph{riduce ogni feature a un indice
categorico}. In uscita ha un oggetto unico, \texttt{LMIndexes}, il cui carico
portante è la \emph{matrice degli indici}: un array di \texttt{uint32} di
forma $(M, I)$, dove $M$ è il numero di posizioni di addestramento e $I$ il
numero totale di istanze di tutte le feature. Ogni \emph{riga} è una
posizione; ogni \emph{colonna} è una sotto-variabile (una istanza di una
feature); ogni \emph{cella} è l'indice categorico assunto da quella
sotto-variabile in quella posizione. È la matrice che, con la giusta
espansione, diventa la matrice di disegno del modello.

\subsection{La matrice degli indici}

La matrice è prodotta da
\texttt{feature\_set.compute\_indexes(positions)}, che impila lungo l'ultimo
asse i blocchi delle feature, in un \emph{layout a colonne} fisso e
documentato:

\begin{itemize}
  \item \textbf{colonna 0} (se presente): l'intercetto, costante \texttt{0};
  \item poi una \emph{colonna per mobility}, nell'ordine del
        \texttt{MobilitySet}: l'indice è il \emph{conteggio} del numero di
        mosse disponibili sul lato scelto, per cui con $k$ caselle nella
        maschera l'indice gira su $0 \ldots k$ (ogni valore raggiungibile);
  \item poi, per ogni pattern dell'ordine per maschera del
        \texttt{PatternSet}, \emph{$n_{\text{inst}}$ colonne}: l'\emph{indice
        principale} di ogni istanza dell'orbita $D_4$, nell'ordine delle istanze
        del pattern.
\end{itemize}

Con un feature set di esempio — intercetto, la mobility \texttt{LMC}, il
pattern \texttt{EDGE} a $8$ istanze e il pattern \texttt{COR} a $4$ istanze —
la matrice ha $I = 1 + 1 + 8 + 4 = 14$ colonne, così disposte:

\begin{center}
\small
\begin{tabular}{l|c|c|c|c}
\toprule
Colonna & 0 & 1 & $2 \ldots 9$ & $10 \ldots 13$ \\
\midrule
Feature & INTERCEPT & LMC & EDGE (ist.\ 0--7) & COR (ist.\ 0--3) \\
\bottomrule
\end{tabular}
\end{center}

Ricordando le definizioni dei capitoli sulla scacchiera e sui pattern: una
\emph{configurazione} di un pattern è l'assegnazione in base 3 delle sue $S$
caselle (0 = vuota, 1 = mover, 2 = avversario), e l'\emph{indice principale}
è il minimo, tra gli $8$ immagini $D_4$ della configurazione, secondo
l'ordinamento che il stabilizzatore della maschera induce sulle caselle.
Il \emph{guadagno} di questa riduzione è che due configurazioni
\emph{equivalenti per simmetria} — anche se su istanze \emph{diverse} del
pattern — condividono lo \emph{stesso} indice principale: è ciò che condensa
al massimo il numero di pesi del modello, come già notato nel capitolo sulla
funzione di valutazione.

\subsection{Lo spazio degli indici}

I due tipi di feature abitano spazi di indice di natura diversa, e la
distinzione sarà cruciale per le mappe di peso.

\emph{Mobility.} Lo spazio è $0 \ldots k$ ($k+1$ valori): il conteggio grezzo,
senza alcuna riduzione per simmetria. La dimensione teorica è
$N_{\text{conf}} = k + 1$.

\emph{Pattern.} Lo spazio in base 3 grezzo ha $3^{S}$ codici, ma gli indici
\emph{principali} occupano solo i primi $n_{\text{unique}}$ di essi (le classi
di equivalenza sotto le simmetrie, $n_{\text{unique}} \le 3^{S}$). L'indice è
un \texttt{Index} (\texttt{uint32}), che resta al sicuro finché $S \le 18$
($3^{18} \approx 3.9 \times 10^{8} < 2^{32}$). I codici in
$[n_{\text{unique}}, 3^{S})$ non vengono \emph{mai} prodotti da un indice
principale: sono, di fatto, irraggiungibili.

Questa è la differenza esatta tra le due grandezze che i livelli analitici
maneggiano: $n_{\text{conf}} = 3^{S}$ (lo spazio in base 3 \emph{grezzo},
usato per \emph{dimensionare} le mappe) e $n_{\text{unique}}$ (le classi di
\emph{equivalenza}, usate per \emph{misurare} la coverage e la rarità di un
pattern).

\subsection{L'oggetto \texttt{LMIndexes} e la cache}

L'oggetto \texttt{LMIndexes} impacchetta la matrice insieme a due \emph{firme}:
il \texttt{rds\_checksum} (ereditato da \texttt{POSITIONS}) e lo
\texttt{feature\_set\_hash} (ereditato dalla fase Setup). La checksum dei dati
è lo SHA3-256 dei \emph{byte grezzi} della matrice. La consistenza della cache
è la congiunzione delle due firme: se \emph{entrambe} la checksum del data set
e lo hash del feature set coincidono con quelle live, la cache è valida. È la
stessa forma di ``chiave composita'' che si è vista per il data set:
garantisce che la matrice cacheata corrisponda \emph{sia} alle stesse
posizioni \emph{sia} allo stesso feature set.

\section{Il livello \texttt{ANOVA}}
\label{sec:corpus-anova}

\texttt{ANOVA} (Analisi della Varianza) è il secondo livello opzionale della
fase. Risponde a una domanda complementare a quella dell'\texttt{SU\_ANALYSIS}:
non ``quale feature duplica quale'', ma \emph{quale feature spiega quanta
varianza del valore di gioco}. È un livello \emph{diagnostico}: non pubblica
nulla sul contesto, scrive solo i propri file di report. Come
l'\texttt{SU\_ANALYSIS}, è opt-in (\texttt{enable} \texttt{false} di default).

L'analisi ha \emph{due passate}, che il worker esegue in sequenza.

\subsection{Le misure di effetto: eta quadrato e omega quadrato}

La \emph{prima} passata è una ANOVA a un fattore \emph{per ogni
sotto-variabile}. Per ogni (feature, istanza) — escluso l'intercetto, che non
è un blocco della matrice — si prende la colonna di indici della matrice come
\emph{assegnazione ai gruppi} e il vettore dei valori di gioco come \emph{target}
continuo. Si partiziona quindi la varianza totale del target in una parte
\emph{spiegata} dall'appartenenza al gruppo (alla configurazione) e una parte
di \emph{residuo}:

\begin{align}
  SS_{\text{tot}} &= \sum_{i} (y_i - \overline{y})^2, \\
  SS_{\text{eff}} &= \sum_{g} n_g\,(\overline{y}_g - \overline{y})^2, \\
  SS_{\text{err}} &= SS_{\text{tot}} - SS_{\text{eff}},
\end{align}

dove la somma efficace è sulle $K$ configurazioni osservate, $n_g$ è il numero
di posizioni nel gruppo $g$ e $\overline{y}_g$ la media del gruppo. Da queste si
ricavano le due \emph{misure di effetto} riportate per ogni sotto-variabile.
L'\emph{eta quadrato}, la frazione della varianza totale spiegata dal fattore
(una stima \emph{con bias}, perché campionaria):

\begin{equation}
  \eta^2 \;=\; \frac{SS_{\text{eff}}}{SS_{\text{tot}}},
  \label{eq:eta2}
\end{equation}

e l'\emph{omega quadrato}, la sua stima \emph{non distorta} (corretta per il
numero di gruppi $K$ e per il residuo medio per grado di libertà
$MS_{\text{err}} = SS_{\text{err}}/(N-K)$):

\begin{equation}
  \omega^2 \;=\; \frac{SS_{\text{eff}} - (K-1)\,MS_{\text{err}}}{SS_{\text{tot}} + MS_{\text{err}}}.
  \label{eq:omega2}
\end{equation}

L'eta quadrato \emph{cresce} al crescere di $K$ anche a parità di informazione
reale (più gruppi = più varianza spiegata per puro caso); l'omega quadrato
paga quel prezzo e può scendere leggermente sotto zero per un effetto nullo —
il valore grezzo è riportato così com'è. Due casi degeneri sono trattati per
convenzione: un solo gruppo osservato ($K = 1$) e una varianza totale nulla
restituiscono $0$ per entrambe le misure. La coppia $(\eta^2, \omega^2)$ per
sotto-variabile è la \emph{metrica di discriminazione} dell'assegnazione di
gruppi: una feature con $\omega^2$ alto è una feature che, da sola, ordina il
valore di gioco. Questa passata \emph{non} ha bisogno del corpus.

\subsection{L'aderenza: DB contro corpus}

La \emph{seconda} passata — l'\emph{aderenza} — è quella che si nutre del
corpus dell'\texttt{SU\_ANALYSIS}. Per ogni sotto-variabile confronta due
distribuzioni categoriche: la \emph{distribuzione osservata} nel set di
addestramento (le frequenze dei codici della colonna, misurate sul data set)
e la \emph{distribuzione ideale} del corpus (le frequenze marginali che
l'\texttt{SU\_ANALYSIS} ha misurato sul campione casuale neutro). Le due sono
allineate sull'\emph{unione} delle loro categorie, e vengono confrontate con
un trio fisso di distanze e test:

\begin{itemize}
  \item la \emph{distanza di variazione totale}, robusta e normalizzata in
        $[0,1]$ (nessun log, nessun termine degenere $0/0$):
        $\mathrm{TV}(p,q) = \tfrac{1}{2}\sum_i |p_i - q_i|$;
  \item la \emph{divergenza di Kullback-Leibler}
        $\mathrm{KL}(p\,\|\,q) = \sum_{i:\,p_i>0} p_i \ln \frac{p_i}{q_i}$,
        in nats, con la convenzione $0\log 0 = 0$ e valore $+\infty$ quando il
        data set mostra un codice che l'ideale non mostra mai;
  \item il \emph{test del chi-quadrato} di bontà di adattamento,
        $\chi^2 = \sum_i (O_i - E_i)^2 / E_i$ con $E_i = p^{\text{ideal}}_i
        \cdot N_{\text{db}}$; le celle la cui frequenza attesa scende sotto una
        soglia minima (default $5$) sono \emph{fuse} in un unico bin ``altro''
        prima del test (la tecnica standard per le celle a bassa attesa), e il
        valore-$p$ è valutato con $k-1$ gradi di libertà.
\end{itemize}

Le sotto-variabili sono poi \emph{ordinate per TV decrescente} e il report
riporta le prime \texttt{top\_k\_adherence} (default 20): sono le feature la
cui distribuzione nel set di addestramento \emph{si discosta} di più dalla
distribuzione naturale del gioco. Una feature in cima a questa classifica è un
segnale d'allarme: il set di addestramento la sta campionando in modo
\emph{non} rappresentativo (per eccesso o per difetto) rispetto al gioco
casuale, e i pesi che il modello imparerà da quelle configurazioni saranno
condizionati da quella selezione.

\subsection{Degradazione elegante e report}

Qui vive l'accorgimento di \emph{degradazione elegante} preannunciato
all'inizio. Se l'\texttt{SU\_ANALYSIS} era \emph{spento} (quindi
\texttt{ctx.su\_analysis\_tables} è \texttt{None} e non c'è corpus nel
contesto), la passata di aderenza è \emph{saltata} con una nota nel log, e
l'ANOVA \emph{prosegue} con la sola passata eta/omega — che, come visto, non
del corpus si serve. Il report TXT (più un CSV opzionale) contiene la sezione
eta/omega per tutte le sotto-variabili e, quando presente, la classifica di
aderenza; quando assente, la riga esplicita
``Adherence: SKIPPED''. In nessun caso il livello fallisce per la mancanza del
corpus.

\section{Il livello \texttt{WMAPS}}
\label{sec:corpus-wmaps}

\texttt{WMAPS} (Weight Maps) è l'ultimo dei tre livelli obbligatori, e il
ponte \emph{esatto} verso la fase di costruzione del modello. Il suo compito è
\emph{costruire le mappe che collegano ogni configurazione di feature a un
indice del vettore di pesi} $\bm{w}$. Ricordando il modello a pattern del
capitolo sulla funzione di valutazione: il predittore lineare somma, per ogni
feature attiva, il peso della sua configurazione corrente. Le mappe di peso
sono il \emph{dizionario} che, data una configurazione, restituisce
\emph{dove} andare a leggere quel peso nel vettore $\bm{w}$.

\subsection{A cosa servono le mappe di peso}

Il problema che le mappe risolvono è uno di \emph{indirizzamento} sparso.
Le colonne della matrice degli indici contengono codici che vivono in spazi
\emph{disgiunti} e di dimensioni diverse (l'intercetto in $\{0\}$, una mobility
in $0\ldots k$, un pattern in $0\ldots 3^{S}-1$), ma il vettore di pesi
$\bm{w}$ è un \emph{unico} array lineare. Serve quindi un \emph{offset} per
feature che sposti lo spazio dei codici di ogni feature nella sua \emph{fascia}
privata del vettore, e una \emph{mappa} che, dentro la fascia, colleghi ogni
configurazione all'effettiva posizione del peso — con la possibilità di
\emph{fare saltare} le configurazioni rare (affidandole a un peso di
\emph{fallback} condiviso) per non allocare un peso proprio a codici che il
data set mostra pochissimo o per nulla.

\subsection{I sei array}

L'oggetto \texttt{LMWMaps} porta \emph{sei} array (più due scalari di
metadati: la checksum degli indici e il \texttt{cut\_off} usato). Con
$F$ il numero di feature, $K = \sum_f n_{\text{conf},f}$ il totale delle
configurazioni e $W$ il totale dei pesi allocati:

\begin{center}
\small
\begin{tabularx}{\textwidth}{@{}llX@{}}
\toprule
Array & Forma & Contenuto \\
\midrule
\texttt{feature\_w\_ranges} & $(F, 3)$ & per feature, in ordine di \texttt{fid}: $[\text{fallback}, w_{\min}, w_{\max}]$, gli indici di peso della fascia. \\
\texttt{iwmap\_feature\_offset} & $(F+1,)$ & la posizione, nell'array \texttt{iwmap}, della configurazione zero della feature \texttt{fid} (somme cumulative delle dimensioni). \\
\texttt{iwmap} & $(K,)$ & per ogni configurazione, l'indice di peso; $-1$ se la configurazione non è presente nel data set. \\
\texttt{wmap} & $(W, 3)$ & per ogni indice di peso: $[\text{fid}, \text{config\_id} \text{ (o } -1 \text{)}, \text{frequenza}]$. \\
\texttt{wmap\_fallback} & $(F_{\text{cut}}, 3)$ & le configurazioni \emph{escluse} dal cut-off, una riga per ciascuna: $[\text{fid}, \text{config\_id}, \text{frequenza}]$. \\
\texttt{w} & $(W,)$ & il vettore dei pesi del modello, \emph{inizializzato a zero}. \\
\bottomrule
\end{tabularx}
\end{center}

Il percorso di lettura è in due passi. Dato un \texttt{fid} e un codice di
configurazione $c$ di quella feature, l'indice di peso si ottiene da

\begin{equation}
  w \;=\; \texttt{iwmap}\big[\, \texttt{iwmap\_feature\_offset}[\text{fid}] + c \,\big].
  \label{eq:iwmap-lookup}
\end{equation}

L'array \texttt{iwmap} è dimensionato su $K$, il totale \emph{grezzo} delle
configurazioni ($\sum_f n_{\text{conf},f}$): per un pattern occupa quindi
l'intero spazio in base 3 $3^{S}$, anche se gli indici principali
\emph{vissuti} stanno solo nei primi $n_{\text{unique}}$ slot. La scelta è
deliberata: dimensionare la fascia sullo spazio grezzo rende la
\eqref{eq:iwmap-lookup} un \emph{accesso diretto} $O(1)$ (una sola somma),
senza tabelle di ricerca; gli slot irraggiungibili
($c \ge n_{\text{unique}}$) restano semplicemente alla sentinella $-1$, che
indica ``configurazione non vista''.

I valori \emph{negativi} che compaiono in \texttt{iwmap} e
\texttt{feature\_w\_ranges} sono \texttt{ExtendedIndex} (\texttt{int64}): qui
$-1$ è la \emph{sentinella logica} ``configurazione non presente nel data
set'', \emph{non} un dato. È l'unico posto della pipeline in cui lo spazio
esteso degli indici (con i negativi) viene effettivamente popolato: la matrice
degli indici in ingresso è \emph{densa} per contratto (ogni cella è un codice
valido in $0 \ldots n_{\text{conf}}-1$), quindi la sentinella può comparire
\emph{solo} in uscita, nelle mappe.

\subsection{Il cut-off di frequenza e il fallback}

Il motore di allocazione è governato dalla soglia
\texttt{stat\_model.frequency\_cut\_off} (default $1$). Per ogni feature,
l'algoritmo \emph{poola} le colonne di \emph{tutte} le sue istanze (le
schiaccia in un unico vettore e conta le frequenze dei codici): è corretto
proprio \emph{perché} le istanze di una feature condividono lo \emph{stesso}
spazio di indici principali, come si è visto al livello
\texttt{INDEXES} — una configurazione che appare sull'istanza 2 e una
configurazione $D_4$-equivalente sull'istanza 5 sono lo \emph{stesso} codice e
concorrono alla \emph{stessa} frequenza.

Sulla base delle frequenze, le configurazioni si dividono in due classi:

\begin{itemize}
  \item le configurazioni con frequenza $< \text{cut\_off}$ (le \emph{rare}):
        \emph{condividono} l'\emph{unico} peso di \emph{fallback} della
        feature. Quel peso è allocato \emph{per primo} nella fascia, la sua
        riga in \texttt{wmap} porta $\text{config\_id} = -1$ e la
        \emph{somma} delle frequenze di tutte le configurazioni rare, e
        ognuna di loro è inoltre elencata in \texttt{wmap\_fallback}. Il
        campo \texttt{fallback} di \texttt{feature\_w\_ranges} ne memorizza
        l'indice (e vale $-1$ se la feature non ha configurazioni rare);
  \item le configurazioni con frequenza $\ge \text{cut\_off}$ (le
        \emph{comuni}): ognuna riceve un indice di peso \emph{contiguo}
        proprio, allocato in blocco subito dopo il (eventuale) fallback.
\end{itemize}

L'effetto complessivo è una \emph{compressione sparsa} del vettore di pesi:
i pesi totali $W$ sono \emph{minori} o uguali al totale grezzo $K$ di
configurazioni, perché ogni feature con configurazioni rare le
\emph{collassa} in un singolo peso di fallback. Al momento dell'\emph{inference},
quando il modello incontra una configurazione non vista in addestramento
($\texttt{iwmap} = -1$), la lettura \eqref{eq:iwmap-lookup} fallisce e si
riporta al \texttt{fallback} memorizzato in
\texttt{feature\_w\_ranges}: la configurazione rara o inedita eredita il peso
appreso sull'\emph{insieme} delle configurazioni rare della sua feature,
invece di contribuire con un peso nullo e non appreso. È il meccanismo che
mantiene il modello \emph{definito} su ogni configurazione, anche fuori
dal supporto del data set.

Per fissare le idee, un esempio in piccolo. Prendiamo un feature set di tre
feature: l'intercetto, la mobility \texttt{LMC} a maschera piena ($65$
valori, da $0$ a $64$) e il pattern \texttt{COR} a una casella d'angolo
($3$ configurazioni, $4$ istanze — i quattro angoli). Il totale delle
configurazioni è $K = 1 + 65 + 3 = 69$, e gli offset sono
\texttt{iwmap\_feature\_offset} $= [0, 1, 66, 69]$. Con il
\texttt{cut\_off} di default ($1$) ogni configurazione osservata si prende un
peso proprio: $W = 69$, nessuna riga di fallback, e \texttt{iwmap} è la
mappa identità dentro ogni fascia.

Alzi ora la soglia a $200$ occorrenze. Nel set di addestramento supponiamo
che i valori estremi della \texttt{LMC} — $0$, $1$, $63$ e $64$ mosse —
siano quelli rari (frequenza $< 200$), mentre gli altri $61$ valori
superino la soglia; le tre configurazioni del \texttt{COR} superino tutte la
soglia. L'allocazione diventa:

\begin{itemize}
  \item \texttt{feature\_w\_ranges}: l'intercetto $[-1, 0, 0]$; la
        \texttt{LMC} $[1, 1, 62]$ — il fallback sta a $1$, la prima
        posizione della fascia; il \texttt{COR} $[-1, 63, 65]$;
  \item \texttt{iwmap}: nella fascia \texttt{LMC} (offset $1$) i quattro
        codici rari puntano tutti al fallback,
        $\texttt{iwmap}[1] = \texttt{iwmap}[2] = \texttt{iwmap}[64] =
        \texttt{iwmap}[65] = 1$, mentre i codici $2 \ldots 62$ puntano a
        $2 \ldots 62$;
  \item \texttt{wmap}: la riga del fallback è
        $[\text{LMC}, -1, f_0 + f_1 + f_{63} + f_{64}]$, con la \emph{somma}
        delle frequenze rare;
  \item \texttt{wmap\_fallback}: le quattro righe
        $[\text{LMC}, 0, f_0]$, $[\text{LMC}, 1, f_1]$,
        $[\text{LMC}, 63, f_{63}]$, $[\text{LMC}, 64, f_{64}]$ — il registro
        di chi ha ereditato il fallback.
\end{itemize}

Il vettore di pesi scende da $69$ a $W = 1 + 62 + 3 = 66$: tre pesi in
meno, scambiati con l'unica riga di fallback. A \emph{inference}, se il
modello incontra uno di quei valori rari della \texttt{LMC} — o una
configurazione del \texttt{COR} mai vista in addestramento — la lettura
\eqref{eq:iwmap-lookup} restituisce la sentinella $-1$ e si ripiega sul
\texttt{fallback} di fascia: il contributo al predittore lineare c'è
sempre, e porta con sé l'esperienza \emph{media} delle configurazioni rare
della feature.

Il vettore \texttt{w} esce da questo livello \emph{tutto a zero}:
\texttt{WMAPS} ne fissa solo la \emph{dimensione} $W$ e la
\emph{topologia} (chi sta dove), lasciando il \emph{valore} dei pesi alla
fase di costruzione del modello, che li ottimizzerà. La checksum dell'oggetto
è lo SHA3-256 dei sei array in ordine di dichiarazione (i due scalari di
metadati non sono hashati); la consistenza della cache è la congiunzione della
checksum degli indici in ingresso e del \texttt{cut\_off}.

\section{Il passaggio alla fase di costruzione del modello}

A fase conclusa, il contesto porta il pacchetto completo di dati
\emph{pronti} per la costruzione del modello. Le tre leve
\emph{obbligatorie} ci sono tutte: \texttt{positions} e
\texttt{game\_values} (il \emph{cosa} addestrare e il \emph{target}),
\texttt{hlp\_indexes} (le feature ridotte a indici categorici) e le sei mappe
di \texttt{WMAPS} (la topologia del vettore di pesi, con $\texttt{w}$ a
zero). Le due leve \emph{opzionali} — le tabelle di
\texttt{SU\_ANALYSIS} e i report di \texttt{ANOVA} — sono a disposizione
come \emph{diagnostica} per giudicare la scelta del feature set, ma non
entrano nel modello.

Il livello immediatamente successivo, \texttt{DESIGN\_MATRIX} (il primo della
fase \emph{Costruzione del modello}), consumerà esattamente
\texttt{hlp\_indexes} e le mappe di peso per \emph{espandere} la matrice
compatta degli indici nella matrice di disegno binaria a $W$ colonne, e
\texttt{game\_values} per costruire il target. In altre parole: questa fase ha
\emph{preparato il terreno e i mattoni}; la fase successiva li impilerà. La
\eqref{eq:iwmap-lookup} che qui è stata introdotta come \emph{lettura} delle
mappe diventerà, nel costruttore del modello, la regola con cui ogni riga
della matrice di disegno sa \emph{quale} colonna di peso attivare.
